\chapter{Boolean Analysis — Restriction}
%%% generated by the blueprint pipeline from TCSlib.BooleanAnalysis.Switching.Restriction
\section{Overview}

This module sets up the basic machinery for the switching lemma: restrictions of
Boolean variables, their effect on literals, terms and DNF formulas, the restricted
Boolean function, the notion of a bad restriction, and a collection of small auxiliary
lemmas about list operations and decision-tree depth used in the main argument.

%%%%%%%%%%%%%%%%%%%%%%%%%%%%%%%%%%%%%%%%%%%%%%%%%%%%%%%%%%%%%%%%%%%%%%%%%%%%%%%
\section{Declarations}

\begin{definition}[Restriction]
\lean{SwitchingLemma2.Restriction}
\statementsource{boolean-ch04-dnf-switching-lmn}{pdf-fc447f326910-p006-b003}
\label{SwitchingLemma2.Restriction}
\leanok
\statementsource{switching-lemma}{pdf-b5e074215b9e-p001-b007}
A restriction on $n$ variables is a map $\rho : \mathrm{Fin}\,n \to \mathrm{Option}\,\mathrm{Bool}$,
assigning each variable either a fixed Boolean value or leaving it free (\texttt{none}).
\end{definition}

\begin{definition}[Free variables of a restriction]
\lean{SwitchingLemma2.Restriction.freeVars}
\statementsource{boolean-ch04-dnf-switching-lmn}{pdf-fc447f326910-p006-b003}
\label{SwitchingLemma2.Restriction.freeVars}
\leanok
\uses{SwitchingLemma2.Restriction}
The set of free variables of a restriction $\rho$ is the finite set of indices
$i \in \mathrm{Fin}\,n$ for which $\rho(i)$ is unset (\texttt{none}).
\end{definition}

\begin{definition}[Number of free variables]
\lean{SwitchingLemma2.Restriction.numFree}
\label{SwitchingLemma2.Restriction.numFree}
\leanok
\uses{SwitchingLemma2.Restriction, SwitchingLemma2.Restriction.freeVars}
The number of free variables of a restriction $\rho$ is the cardinality of its set of
free variables.
\end{definition}

\begin{definition}[Extension of a point along a restriction]
\lean{SwitchingLemma2.Restriction.extend}
\label{SwitchingLemma2.Restriction.extend}
\leanok
\uses{SwitchingLemma2.Restriction}
Given a restriction $\rho$ and a point $x : \mathrm{Fin}\,n \to \mathrm{Bool}$, the
extension $\rho.\mathrm{extend}\,x$ is the Boolean assignment that uses the fixed value
$\rho(i)$ where it is set and falls back to $x(i)$ on the free variables.
\end{definition}

\begin{definition}[Fixing variables in a restriction]
\lean{SwitchingLemma2.Restriction.fixVars}
\label{SwitchingLemma2.Restriction.fixVars}
\leanok
\uses{SwitchingLemma2.Restriction}
Given a list of (variable, value) pairs, $\mathrm{fixVars}$ updates the restriction so
that each listed variable is set to the corresponding fixed Boolean value.
\end{definition}

\begin{definition}[Un-fixing variables in a restriction]
\lean{SwitchingLemma2.Restriction.unfixVars}
\label{SwitchingLemma2.Restriction.unfixVars}
\leanok
\uses{SwitchingLemma2.Restriction}
Given a list of (variable, value) pairs, $\mathrm{unfixVars}$ resets each listed variable
to \texttt{none}, restoring it to a free variable.
\end{definition}

\begin{definition}[Restriction with $s$ free variables]
\lean{SwitchingLemma2.IsRestriction}
\label{SwitchingLemma2.IsRestriction}
\leanok
\statementsource{switching-lemma}{pdf-b5e074215b9e-p001-b007}
\uses{SwitchingLemma2.Restriction, SwitchingLemma2.Restriction.numFree}
The predicate $\mathrm{IsRestriction}\,s\,\rho$ holds when the restriction $\rho$ has
exactly $s$ free variables.
\end{definition}

\begin{definition}[Literal killed by a restriction]
\lean{SwitchingLemma2.Literal.killedBy}
\label{SwitchingLemma2.Literal.killedBy}
\leanok
\statementsource{switching-lemma}{pdf-b5e074215b9e-p002-b002}
\uses{Literal, SwitchingLemma2.Restriction}
A literal $l$ is killed by $\rho$ when $\rho$ fixes its variable to the value that makes
$l$ false; concretely $\rho(l.\mathrm{var}) = \mathrm{some}\,l.\mathrm{neg}$.
\end{definition}

\begin{definition}[Literal fixed by a restriction]
\lean{SwitchingLemma2.Literal.fixedBy}
\label{SwitchingLemma2.Literal.fixedBy}
\leanok
\statementsource{switching-lemma}{pdf-b5e074215b9e-p002-b002}
\uses{Literal, SwitchingLemma2.Restriction}
A literal $l$ is fixed by $\rho$ when $\rho$ fixes its variable to the value that makes
$l$ true; concretely $\rho(l.\mathrm{var}) = \mathrm{some}\,(\lnot\, l.\mathrm{neg})$.
\end{definition}

\begin{definition}[Term killed by a restriction]
\lean{SwitchingLemma2.Term.killedBy}
\label{SwitchingLemma2.Term.killedBy}
\leanok
\statementsource{switching-lemma}{pdf-b5e074215b9e-p002-b002}
\uses{SwitchingLemma2.Literal.killedBy, SwitchingLemma2.Restriction, Term}
A term (conjunction of literals) is killed by $\rho$ when at least one of its literals is
killed by $\rho$, forcing the whole term to evaluate to false.
\end{definition}

\begin{definition}[Term fixed by a restriction]
\lean{SwitchingLemma2.Term.fixedBy}
\label{SwitchingLemma2.Term.fixedBy}
\leanok
\statementsource{switching-lemma}{pdf-b5e074215b9e-p002-b002}
\uses{SwitchingLemma2.Literal.fixedBy, SwitchingLemma2.Restriction, Term}
A term is fixed by $\rho$ when every one of its literals is fixed by $\rho$, forcing the
whole term to evaluate to true.
\end{definition}

\begin{definition}[Alive term]
\lean{SwitchingLemma2.isAlive}
\label{SwitchingLemma2.isAlive}
\leanok
\statementsource{switching-lemma}{pdf-b5e074215b9e-p002-b002}
\uses{SwitchingLemma2.Restriction, SwitchingLemma2.Term.fixedBy, SwitchingLemma2.Term.killedBy, Term}
A term $t$ is alive under $\rho$ when it is neither killed nor fixed by $\rho$.
\end{definition}

\begin{definition}[Restricted Boolean function]
\lean{SwitchingLemma2.restrictFn}
\statementsource{boolean-ch04-dnf-switching-lmn}{pdf-fc447f326910-p006-b004}
\label{SwitchingLemma2.restrictFn}
\leanok
\statementsource{switching-lemma}{pdf-b5e074215b9e-p001-b007}
\uses{SwitchingLemma2.Restriction, SwitchingLemma2.Restriction.extend}
Given $f : (\mathrm{Fin}\,n \to \mathrm{Bool}) \to \mathrm{Bool}$ and a restriction
$\rho$, the restricted function maps $x$ to $f(\rho.\mathrm{extend}\,x)$, i.e.\ $f$
evaluated on the free coordinates with the fixed coordinates set by $\rho$.
\end{definition}

\begin{definition}[Bad restriction]
\lean{SwitchingLemma2.IsBadRestriction}
\statementsource{boolean-ch04-dnf-switching-lmn}{
  pdf-fc447f326910-p008-b005,
  pdf-fc447f326910-p018-b002
}
\label{SwitchingLemma2.IsBadRestriction}
\leanok
\statementsource{switching-lemma}{pdf-b5e074215b9e-p002-b001}
\uses{SwitchingLemma2.Restriction, SwitchingLemma2.restrictFn, dtDepth}
A restriction $\rho$ is bad for $f$ at depth $d$ when the restricted function
$\mathrm{restrictFn}\,f\,\rho$ has decision-tree depth strictly greater than $d$.
\end{definition}

\begin{definition}[Number of $s$-restrictions]
\lean{SwitchingLemma2.numSRestrictions}
\label{SwitchingLemma2.numSRestrictions}
\leanok
The number of restrictions on $n$ variables leaving exactly $s$ free, namely
$\binom{n}{s}\, 2^{\,n-s}$.
\end{definition}

\begin{lemma}[Killed literals vanish under the induced assignment]
\lean{SwitchingLemma2.Literal.killedBy_eval_false}
\label{SwitchingLemma2.Literal.killedBy_eval_false}
\leanok
\uses{Literal, Literal.eval, SwitchingLemma2.Literal.killedBy, SwitchingLemma2.Restriction, SwitchingLemma2.Restriction.extend}
\difficulty{2}
Fix $n$ and let $l$ be a literal on $n$ variables, with underlying variable $v \in
\mathrm{Fin}\,n$ and polarity flag $b \in \{\text{false}, \text{true}\}$ (with $b =
\text{true}$ marking a negated literal). Let $\rho$ be a restriction on $n$ variables
that kills $l$, meaning $\rho$ fixes $v$ to the value $b$. Then for every assignment $x
: \mathrm{Fin}\,n \to \{\text{false}, \text{true}\}$, the literal $l$ evaluates to false
at the extension of $x$ along $\rho$.
\end{lemma}

\begin{lemma}[A fixed literal evaluates to true on any extension]
\lean{SwitchingLemma2.Literal.fixedBy_eval_true}
\label{SwitchingLemma2.Literal.fixedBy_eval_true}
\leanok
\uses{Literal, Literal.eval, SwitchingLemma2.Literal.fixedBy, SwitchingLemma2.Restriction, SwitchingLemma2.Restriction.extend}
\difficulty{2}
Let $l$ be a literal on $n$ variables and let $\rho$ be a restriction on those same
variables that fixes $l$, meaning $\rho$ assigns $l$'s variable the Boolean value that
makes $l$ true. Then for every assignment $x : \mathrm{Fin}\,n \to \mathrm{Bool}$, the
literal $l$ evaluates to true on the extension of $x$ along $\rho$ (the assignment using
$\rho$'s fixed value on each set variable and $x$ on the free ones).
\end{lemma}

\begin{lemma}[Decision-tree depth bound from a computing tree]
\lean{SwitchingLemma2.dtDepth_le_of_tree}
\label{SwitchingLemma2.dtDepth_le_of_tree}
\leanok
\uses{DecisionTree, DecisionTree.depth, DecisionTree.eval, dtDepth}
\difficulty{2}
Let $f : (\mathrm{Fin}\,n \to \mathrm{Bool}) \to \mathrm{Bool}$ be a Boolean function,
and let $T$ be a decision tree on $n$ variables of depth at most $d$ that computes $f$,
meaning $T$ evaluates to $f(x)$ on every input $x$. Then the decision-tree depth of $f$
is at most $d$.
\end{lemma}

\begin{lemma}[Boolean existential over a list is false when the predicate always fails]
\lean{SwitchingLemma2.list_any_eq_false}
\label{SwitchingLemma2.list_any_eq_false}
\leanok
\difficulty{3}
Let $l$ be a finite list of elements of a type $\alpha$, and let $p \colon \alpha \to
\{\mathrm{true}, \mathrm{false}\}$ be a Boolean predicate. If $p(x) = \mathrm{false}$
for every entry $x$ of $l$, then the Boolean existential $\bigvee_{x \in l} p(x)$ — the
value $\mathrm{true}$ exactly when some entry of $l$ satisfies $p$ — is
$\mathrm{false}$.
\end{lemma}

\begin{lemma}[A false witness makes a universal test fail]
\lean{SwitchingLemma2.list_all_eq_false_of_mem}
\label{SwitchingLemma2.list_all_eq_false_of_mem}
\leanok
\difficulty{4}
Let $l$ be a list of elements of a type $\alpha$ and let $p \colon \alpha \to
\{\mathrm{true}, \mathrm{false}\}$ be a Boolean predicate. If some element $a$ of $l$
satisfies $p(a) = \mathrm{false}$, then the test whether every entry of $l$ satisfies
$p$ evaluates to $\mathrm{false}$.
\end{lemma}

\begin{lemma}[A fixed term forces zero decision-tree depth]
\lean{SwitchingLemma2.fixedTerm_implies_dtDepth_zero}
\label{SwitchingLemma2.fixedTerm_implies_dtDepth_zero}
\leanok
\proofsource{switching-lemma}{pdf-b5e074215b9e-p002-b002}
\proofstep
  {BoolCircuit.Lit}
  {context}
  {switching-lemma}
  {pdf-b5e074215b9e-p001-b004}
\proofstep
  {BoolCircuit.Lit.eval}
  {context}
  {switching-lemma}
  {pdf-b5e074215b9e-p001-b004}
\proofstep
  {Literal}
  {context}
  {switching-lemma}
  {pdf-b5e074215b9e-p001-b004}
\proofstep
  {DecisionTree.eval}
  {context}
  {switching-lemma}
  {pdf-b5e074215b9e-p001-b005}
\proofstep
  {Term}
  {context}
  {switching-lemma}
  {pdf-b5e074215b9e-p001-b004}
\proofstep
  {DNF}
  {context}
  {switching-lemma}
  {pdf-b5e074215b9e-p001-b004}
\proofstep
  {DecisionTree}
  {context}
  {switching-lemma}
  {pdf-b5e074215b9e-p001-b005}
\proofstep
  {DecisionTree.depth}
  {context}
  {switching-lemma}
  {pdf-b5e074215b9e-p001-b005}
\proofstep
  {buildFullDTree}
  {context}
  {switching-lemma}
  {pdf-b5e074215b9e-p001-b005}
\proofstep
  {buildFullDTree_depth}
  {context}
  {switching-lemma}
  {pdf-b5e074215b9e-p001-b005}
\proofstep
  {buildFullDTree_eval}
  {context}
  {switching-lemma}
  {pdf-b5e074215b9e-p001-b005}
\proofstep
  {dtDepth}
  {context}
  {switching-lemma}
  {pdf-b5e074215b9e-p001-b005}
\proofstep
  {SwitchingLemma2.Restriction}
  {context}
  {switching-lemma}
  {pdf-b5e074215b9e-p001-b007}
\proofstep
  {SwitchingLemma2.Restriction.extend}
  {context}
  {switching-lemma}
  {pdf-b5e074215b9e-p001-b007}
\proofstep
  {SwitchingLemma2.Literal.fixedBy}
  {context}
  {switching-lemma}
  {pdf-b5e074215b9e-p001-b007}
\proofstep
  {SwitchingLemma2.Term.fixedBy}
  {context}
  {switching-lemma}
  {pdf-b5e074215b9e-p002-b002}
\proofstep
  {SwitchingLemma2.restrictFn}
  {context}
  {switching-lemma}
  {pdf-b5e074215b9e-p001-b007}
\proofstep
  {SwitchingLemma2.Literal.fixedBy_eval_true}
  {context}
  {switching-lemma}
  {pdf-b5e074215b9e-p002-b002}
\proofstep
  {SwitchingLemma2.dtDepth_le_of_tree}
  {context}
  {switching-lemma}
  {pdf-b5e074215b9e-p001-b005}
\uses{DNF, DNF.eval, DecisionTree.eval, SwitchingLemma2.Literal.fixedBy_eval_true, SwitchingLemma2.Restriction, SwitchingLemma2.Restriction.extend, SwitchingLemma2.Term.fixedBy, SwitchingLemma2.dtDepth_le_of_tree, SwitchingLemma2.restrictFn, Term.eval, dtDepth}
\difficulty{4}
Let $f$ be a DNF formula on $n$ Boolean variables and let $\rho$ be a restriction of
those variables. If at least one term of $f$ is fixed by $\rho$, then the restricted
Boolean function $x \mapsto f(\rho\text{-extension of }x)$ has decision-tree depth $0$.
\end{lemma}

\begin{lemma}[Killed terms force zero decision-tree depth]
\lean{SwitchingLemma2.killedAll_implies_dtDepth_zero}
\label{SwitchingLemma2.killedAll_implies_dtDepth_zero}
\leanok
\proofsource{switching-lemma}{pdf-b5e074215b9e-p002-b002}
\proofstep
  {BoolCircuit.Lit}
  {context}
  {switching-lemma}
  {pdf-b5e074215b9e-p001-b004}
\proofstep
  {BoolCircuit.Lit.eval}
  {context}
  {switching-lemma}
  {pdf-b5e074215b9e-p001-b004}
\proofstep
  {Literal}
  {context}
  {switching-lemma}
  {pdf-b5e074215b9e-p001-b004}
\proofstep
  {DecisionTree.eval}
  {context}
  {switching-lemma}
  {pdf-b5e074215b9e-p001-b005}
\proofstep
  {Term}
  {context}
  {switching-lemma}
  {pdf-b5e074215b9e-p001-b004}
\proofstep
  {DNF}
  {context}
  {switching-lemma}
  {pdf-b5e074215b9e-p001-b004}
\proofstep
  {DecisionTree}
  {context}
  {switching-lemma}
  {pdf-b5e074215b9e-p001-b005}
\proofstep
  {DecisionTree.depth}
  {context}
  {switching-lemma}
  {pdf-b5e074215b9e-p001-b005}
\proofstep
  {buildFullDTree}
  {context}
  {switching-lemma}
  {pdf-b5e074215b9e-p001-b005}
\proofstep
  {buildFullDTree_depth}
  {context}
  {switching-lemma}
  {pdf-b5e074215b9e-p001-b005}
\proofstep
  {buildFullDTree_eval}
  {context}
  {switching-lemma}
  {pdf-b5e074215b9e-p001-b005}
\proofstep
  {dtDepth}
  {context}
  {switching-lemma}
  {pdf-b5e074215b9e-p001-b005}
\proofstep
  {SwitchingLemma2.Restriction}
  {context}
  {switching-lemma}
  {pdf-b5e074215b9e-p001-b007}
\proofstep
  {SwitchingLemma2.Restriction.extend}
  {context}
  {switching-lemma}
  {pdf-b5e074215b9e-p001-b007}
\proofstep
  {SwitchingLemma2.Literal.fixedBy}
  {context}
  {switching-lemma}
  {pdf-b5e074215b9e-p002-b004}
\proofstep
  {SwitchingLemma2.Term.killedBy}
  {context}
  {switching-lemma}
  {pdf-b5e074215b9e-p002-b002}
\proofstep
  {SwitchingLemma2.restrictFn}
  {context}
  {switching-lemma}
  {pdf-b5e074215b9e-p001-b007}
\proofstep
  {SwitchingLemma2.Literal.killedBy_eval_false}
  {context}
  {switching-lemma}
  {pdf-b5e074215b9e-p002-b002}
\proofstep
  {SwitchingLemma2.dtDepth_le_of_tree}
  {context}
  {switching-lemma}
  {pdf-b5e074215b9e-p001-b005}
\proofstep
  {SwitchingLemma2.list_any_eq_false}
  {context}
  {switching-lemma}
  {pdf-b5e074215b9e-p001-b004}
\proofstep
  {SwitchingLemma2.list_all_eq_false_of_mem}
  {context}
  {switching-lemma}
  {pdf-b5e074215b9e-p002-b002}
\uses{DNF, DNF.eval, DecisionTree.eval, SwitchingLemma2.Literal.killedBy_eval_false, SwitchingLemma2.Restriction, SwitchingLemma2.Restriction.extend, SwitchingLemma2.Term.killedBy, SwitchingLemma2.dtDepth_le_of_tree, SwitchingLemma2.list_all_eq_false_of_mem, SwitchingLemma2.list_any_eq_false, SwitchingLemma2.restrictFn, Term.eval, dtDepth}
\difficulty{4}
Let $f$ be a DNF formula on $n$ Boolean variables and let $\rho$ be a restriction on
these variables. If every term of $f$ is killed by $\rho$ — that is, each term contains
a literal whose variable $\rho$ fixes to the value making that literal false — then the
restricted function $x \mapsto f\bigl(\rho.\mathrm{extend}\,x\bigr)$ has decision-tree
depth $0$.
\end{lemma}

\begin{lemma}[Killing preserved under agreement on fixed variables]
\lean{SwitchingLemma2.killedBy_of_nonfree_agree}
\label{SwitchingLemma2.killedBy_of_nonfree_agree}
\leanok
\uses{SwitchingLemma2.Literal.killedBy, SwitchingLemma2.Restriction, SwitchingLemma2.Term.killedBy, Term}
\difficulty{3}
Let $t$ be a term on $n$ variables and let $\rho, \sigma : \mathrm{Fin}\,n \to
\mathrm{Option}\,\mathrm{Bool}$ be two restrictions such that $\sigma(v) = \rho(v)$ for
every variable $v$ that $\rho$ fixes (that is, whenever $\rho(v) \neq \mathrm{none}$).
If $t$ is killed by $\rho$, then $t$ is also killed by $\sigma$.
\end{lemma}

\begin{lemma}[First surviving term is preserved under agreement on fixed variables]
\lean{SwitchingLemma2.first_clause_preserved}
\label{SwitchingLemma2.first_clause_preserved}
\leanok
\uses{DNF, SwitchingLemma2.Restriction, SwitchingLemma2.Term.killedBy, SwitchingLemma2.killedBy_of_nonfree_agree, Term}
\difficulty{4}
Let $f$ be a DNF formula on $n$ variables — a list of terms, each a conjunction of
literals — and let $\rho$ and $\sigma$ be restrictions of these $n$ variables. Suppose
$\sigma$ agrees with $\rho$ on every variable that $\rho$ fixes, that is, $\sigma(v) =
\rho(v)$ for every variable $v$ on which $\rho$ takes a Boolean value rather than
leaving it free. Let $t$ be the first term of $f$ (in list order) that is not killed by
$\rho$, and suppose in addition that $t$ is not killed by $\sigma$. Then $t$ is also the
first term of $f$ that is not killed by $\sigma$.
\end{lemma}

\begin{lemma}[Term length bounded by DNF width]
\lean{SwitchingLemma2.term_length_le_width}
\label{SwitchingLemma2.term_length_le_width}
\leanok
\uses{DNF, DNF.width, Term, Term.width}
\difficulty{3}
Let $f$ be a DNF formula in $n$ variables, that is, a disjunction $t_1 \vee \cdots \vee
t_m$ of terms, where each term is a conjunction of literals, and let its width be the
largest number of literals occurring in any of its terms (and $0$ if $f$ has no terms).
Then for every term $t$ appearing in $f$, the number of literals in $t$ is at most the
width of $f$.
\end{lemma}

\begin{lemma}[Projecting an indexed search back to the list]
\lean{SwitchingLemma2.zipIdx_find_to_find}
\label{SwitchingLemma2.zipIdx_find_to_find}
\leanok
\difficulty{4}
Let $l$ be a list over a type $\alpha$, let
$p\colon\alpha\to\{\text{true},\text{false}\}$ be a predicate, and fix a starting index
$s$. Form the indexed list obtained by pairing each entry of $l$ with its position,
numbered consecutively from $s$. If the first pair in this indexed list whose first
component satisfies $p$ is the pair $(x,\mathrm{idx})$, then $x$ is the first entry of
$l$ satisfying $p$.
\end{lemma}

\begin{lemma}[Membership in the index-tagged list locates a tail]
\lean{SwitchingLemma2.zipIdx_drop_spec}
\label{SwitchingLemma2.zipIdx_drop_spec}
\leanok
\difficulty{3}
Let $t$ be a list over a type $\alpha$, and tag each entry of $t$ with its zero-based
position, forming the list of pairs $(t_i, i)$. If a pair $(l, i)$ occurs among these
tagged entries — equivalently, the entry of $t$ at position $i$ is $l$ — then deleting
the first $i$ entries of $t$ leaves a nonempty list whose first element is $l$; that is,
there is a list $\mathit{rest}$ with $t$ from position $i$ onward equal to $l$ prepended
to $\mathit{rest}$.
\end{lemma}

\begin{lemma}[Position index of a filtered tagged list stays below the length bound]
\lean{SwitchingLemma2.zipIdx_filter_idx_lt}
\label{SwitchingLemma2.zipIdx_filter_idx_lt}
\leanok
\difficulty{3}
Let $t$ be a finite list of elements of a type $\alpha$, and tag each entry of $t$ with
its position to form the list of pairs $(t_i, i)$, whose second coordinates run through
$0, 1, \dots, |t| - 1$. Let $p$ be any Boolean predicate on such pairs, and let $w$ be a
natural number with $|t| \le w$. If a pair $(l, \mathit{idx})$ occurs in the sublist
obtained by keeping exactly those tagged pairs on which $p$ is true, then $\mathit{idx}
< w$.
\end{lemma}
